% !PRINTED 486
\weberpage{486}
%<<CONT
bekommen. Machen wir dann für die sieben Variablen $\eta_r$ die Substitution

\weq{4}{\eta_r = a_0 + a_1 \lambda_r + a_2 \lambda_r^2 + a_3 \lambda_r^3 + a_4 \lambda_r^4 + a_5 \lambda_r^5 + a_6 \lambda_r^6,}

deren Determinante als das Product aller Differenzen $\lambda_i - \lambda_k$ von Null verschieden ist, und substituiren entsprechend

\weq{5}{\eta_r = \sum_{0,\,6}^{s} a_s \lambda_r^s, \quad \eta_r' = \sum_{0,\,6}^{s} a_s' \lambda_r^s, \quad \eta_r'' = \sum_{0,\,6}^{s} a_s'' \lambda_r^s,}

so geht dadurch $X_1$ in eine trilineare Form der drei Variablenreihen $a_s, a_s', a_s''$ über, die nicht identisch verschwinden kann, weil ja auch umgekehrt die $a_s, a_s', a_s''$ durch die $\eta_r, \eta_r', \eta_r''$ linear ausdrückbar sind. Nun kann man für die Variablen $a_s, a_s', a_s''$ solche rationale Zahlenwerthe annehmen (Bd. I, §.\,143), dass die Functionen $X_1, X_2, X_3$ von Null verschiedene Werthe annehmen, und dann stellt (5) eine geeignete Annahme für die Functionen $\eta_r, \eta_r', \eta_r''$ dar\footnote{Vergl. über die hiermit erledigte Frage: \gesperrt{Burckhardt}, \glqq Ueber einen fundamentalen Satz der Lehre von den endlichen Gruppen linearer Substitutionen\grqq{}. Mathem. Annalen, Bd.\ 42 (1892).}.
